\documentclass[11pt,a4paper]{article}
\usepackage[margin=2.3cm]{geometry}
\usepackage{amsmath,amssymb,amsthm,bm}
\usepackage{booktabs}
\usepackage[colorlinks=true,linkcolor=blue,citecolor=blue]{hyperref}

\newcommand{\ii}{\mathrm{i}}
\newcommand{\dd}{\mathrm{d}}
\newcommand{\Res}{\operatorname*{Res}}
\newcommand{\sgn}{\operatorname{sgn}}
\newcommand{\RE}{\operatorname{Re}}
\newcommand{\IM}{\operatorname{Im}}
\newtheorem{remark}{Remark}

\title{Extended Lamb's Problem in the \texorpdfstring{$f$--$k$}{f-k} Domain:\\
Boundary-Condition Systems and Energy Partition (2D and 3D)}
\author{}
\date{\today}

\begin{document}
\maketitle

\tableofcontents

%======================================================================
\section{Conventions}
%======================================================================

\paragraph{Scope.} Sections~\ref{sec:conv-start}--\ref{sec:leaky} treat the 2D P--SV line-source
problem, which gives the building blocks. Section~\ref{sec:3d} extends the formulation to 3D with a
point force of arbitrary direction and an arbitrary moment tensor (any $3\times3$ $\mathbf M$; symmetric in the seismic case). It
reuses the 2D matrices unchanged.

\paragraph{Geometry.}\label{sec:conv-start} Two-dimensional P--SV (plane strain) problem with a line source
parallel to $y$. The interface is $z=0$, with $z$ positive downward. The \emph{solid}
(medium~1: $\rho_1,\alpha_1,\beta_1$, $\mu_1=\rho_1\beta_1^2$, $\lambda_1=\rho_1\alpha_1^2-2\mu_1$)
always occupies $z>0$. The upper half-space $z<0$ is either solid~2
($\rho_2,\alpha_2,\beta_2$) or an inviscid fluid ($\rho_f,\alpha_f$, $\lambda_f=\rho_f\alpha_f^2$).
The source is at $(x,z)=(0,z_s)$, with $z_s=h>0$ (Cases~1--2) or $z_s=-h<0$ (Case~3).

\paragraph{Transforms.} Time dependence $e^{-\ii\omega t}$ (suppressed). Horizontal Fourier transform
\begin{equation}
  \tilde f(k,z)=\int_{-\infty}^{\infty} f(x,z)\,e^{-\ii kx}\,\dd x,
  \qquad
  f(x,z)=\frac{1}{2\pi}\int_{-\infty}^{\infty}\tilde f(k,z)\,e^{\ii kx}\,\dd k ,
\end{equation}
so that $\partial_x\to\ii k$. Tildes are dropped from here on.

\paragraph{Potentials.} $\mathbf u=\nabla\phi+\nabla\times(\psi\,\hat{\mathbf y})$, i.e.
\begin{equation}
  u_x=\partial_x\phi-\partial_z\psi,\qquad u_z=\partial_z\phi+\partial_x\psi,
\end{equation}
with $(\nabla^2+k_\alpha^2)\phi=0$, $(\nabla^2+k_\beta^2)\psi=0$ away from the source, and
$k_m=\omega/c_m$ for $c_m\in\{\alpha,\beta,\alpha_f\}$.

\paragraph{Vertical wavenumbers (Riemann sheet).}
\begin{equation}
  \nu_m(k)=\sqrt{k^2-k_m^2},\qquad \RE\nu_m\ge 0,\quad \IM\nu_m\le 0 ,
  \label{eq:branch}
\end{equation}
so $e^{-\nu_m z}$ is outgoing/decaying for $z\to+\infty$ and $e^{+\nu_m z}$ for $z\to-\infty$.
For $|k|<k_m$: $\nu_m=-\ii\gamma_m$, $\gamma_m=\sqrt{k_m^2-k^2}>0$.
In MATLAB this branch is obtained with \texttt{nu = -1i*sqrt(km\^{}2 - k.\^{}2)}.
We also use
\begin{equation}
  \Gamma_j(k)=2k^2-k_{\beta j}^2,\qquad
  R(k)=\Gamma_1^2-4k^2\nu_{\alpha1}\nu_{\beta1}\quad\text{(Rayleigh function)}.
\end{equation}

\paragraph{Plane-wave fields in the $f$--$k$ domain.}
Using $\lambda k_\alpha^2=\rho\omega^2-2\mu k_\alpha^2$ and $\nu_\alpha^2=k^2-k_\alpha^2$,
\begin{align}
  \sigma_{zz}&=\mu\,\Gamma\,\phi+2\ii\mu k\,\partial_z\psi, &
  \sigma_{xz}&=\mu\bigl(2\ii k\,\partial_z\phi-\Gamma\,\psi\bigr), &
  \sigma_{xx}&=-\lambda k_\alpha^2\phi+2\ii\mu k\,u_x .
\end{align}
For $\phi=a\,e^{s\nu_\alpha z}$, $\psi=b\,e^{s\nu_\beta z}$ with $s=\mp1$ (down-/up-going):

\begin{center}
\renewcommand{\arraystretch}{1.3}
\begin{tabular}{@{}lcc@{}}
\toprule
 & $s=-1$ (field $\propto e^{-\nu z}$) & $s=+1$ (field $\propto e^{+\nu z}$)\\
\midrule
$u_x$ & $\ii k\,a+\nu_\beta b$ & $\ii k\,a-\nu_\beta b$\\
$u_z$ & $-\nu_\alpha a+\ii k\,b$ & $\nu_\alpha a+\ii k\,b$\\
$\sigma_{zz}$ & $\mu\Gamma a-2\ii\mu k\nu_\beta b$ & $\mu\Gamma a+2\ii\mu k\nu_\beta b$\\
$\sigma_{xz}$ & $-2\ii\mu k\nu_\alpha a-\mu\Gamma b$ & $2\ii\mu k\nu_\alpha a-\mu\Gamma b$\\
\bottomrule
\end{tabular}
\end{center}

\noindent Fluid ($\mu=0$, $\psi=0$): $u_z=s\,\nu_f\,a$, $\sigma_{zz}=-p=-\rho_f\omega^2 a$,
$\sigma_{xz}=0$.

%======================================================================
\section{Source representations}
%======================================================================

The 2D Green's function $G$ of $(\nabla^2+k_m^2)G=-\delta(x)\delta(z-z_s)$ is
\begin{equation}
  G(x,z)=\tfrac{\ii}{4}H_0^{(1)}(k_m r)
  \quad\Longleftrightarrow\quad
  \tilde G(k,z)=\frac{e^{-\nu_m|z-z_s|}}{2\nu_m}.
\end{equation}

\paragraph{Explosive (P) line source in a solid or pressure source in a fluid} of strength $\Phi_0$:
\begin{equation}
  \phi^0=\Phi_0\,\frac{e^{-\nu_\alpha|z-z_s|}}{2\nu_\alpha},\qquad \psi^0=0 .
\end{equation}

\paragraph{Vertical line force $F\hat{\mathbf z}$ in a solid (classical Lamb load).}
Writing $F\hat{\mathbf z}\,\delta=\nabla(-F\partial_z g)+\nabla\times(-F\partial_x g\,\hat{\mathbf y})$,
with $\nabla^2 g=-\delta$, the static parts cancel in $\mathbf u$ for $z\ne z_s$, and
\begin{equation}
  \phi^0=\frac{F}{2\rho\omega^2}\,\sgn(z-z_s)\,e^{-\nu_\alpha|z-z_s|},\qquad
  \psi^0=-\frac{\ii kF}{2\rho\omega^2\nu_\beta}\,e^{-\nu_\beta|z-z_s|}.
\end{equation}

\paragraph{Incident amplitudes at the interface.}
For a source in the solid ($z_s=h$) the field travelling toward $z=0$ is
$\phi^0=P^0e^{+\nu_\alpha z}$, $\psi^0=S^0e^{+\nu_\beta z}$ for $0\le z<h$; the field leaving
the source downward is $\phi^0=P_d^0e^{-\nu_\alpha z}$, $\psi^0=S_d^0e^{-\nu_\beta z}$ for $z>h$:
\begin{center}
\renewcommand{\arraystretch}{1.5}
\begin{tabular}{@{}lcccc@{}}
\toprule
Source & $P^0$ & $S^0$ & $P^0_d$ & $S^0_d$\\
\midrule
Explosion & $\dfrac{\Phi_0e^{-\nu_\alpha h}}{2\nu_\alpha}$ & $0$
          & $\dfrac{\Phi_0e^{\nu_\alpha h}}{2\nu_\alpha}$ & $0$\\
Vertical force & $-\dfrac{Fe^{-\nu_\alpha h}}{2\rho\omega^2}$ & $-\dfrac{\ii kFe^{-\nu_\beta h}}{2\rho\omega^2\nu_\beta}$
               & $\dfrac{Fe^{\nu_\alpha h}}{2\rho\omega^2}$ & $-\dfrac{\ii kFe^{\nu_\beta h}}{2\rho\omega^2\nu_\beta}$\\
\bottomrule
\end{tabular}
\end{center}
For a pressure source in the fluid ($z_s=-h$): toward the interface
$\phi^0=P_f^0e^{-\nu_f z}$ with $P_f^0=\Phi_0e^{-\nu_f h}/(2\nu_f)$ ($-h<z\le0$); away from it
$\phi^0=U_f^0e^{+\nu_f z}$ with $U_f^0=\Phi_0e^{\nu_f h}/(2\nu_f)$ ($z<-h$).

\begin{remark}
$e^{+\nu h}$ grows for evanescent $k$, but $P_d^0,S_d^0,U_f^0$ are used only in flux integrals
restricted to $|k|<k_m$, where $|e^{\nu h}|=1$. Never evaluate them outside that window.
\end{remark}

%======================================================================
\section{Boundary-condition systems \texorpdfstring{$\mathbf M(k,\omega)\,\mathbf x=\mathbf S$}{M x = S}}
%======================================================================

The scattered fields are
\begin{align}
  z>0\ \text{(solid 1)}:&\quad \phi_1=A_1e^{-\nu_{\alpha1}z},\quad \psi_1=B_1e^{-\nu_{\beta1}z},\\
  z<0\ \text{(solid 2)}:&\quad \phi_2=A_2e^{+\nu_{\alpha2}z},\quad \psi_2=B_2e^{+\nu_{\beta2}z},\\
  z<0\ \text{(fluid)}:&\quad \phi_f=A_fe^{+\nu_{f}z}.
\end{align}
In every case $\mathbf M$ depends only on the interface; the source enters through $\mathbf S$ only.

\subsection{Case 1: solid--solid, welded contact, source in solid 1}

Continuity of $u_x,u_z,\sigma_{zz},\sigma_{xz}$ at $z=0$ with the total field in medium~1 being
incident plus reflected. Unknowns $\mathbf x=[A_1,\,B_1,\,A_2,\,B_2]^{\mathsf T}$ (reflected P, SV;
transmitted P, SV):
\begin{equation}
\underbrace{\begin{bmatrix}
 \ii k & \nu_{\beta1} & -\ii k & \nu_{\beta2}\\[2pt]
 -\nu_{\alpha1} & \ii k & -\nu_{\alpha2} & -\ii k\\[2pt]
 \mu_1\Gamma_1 & -2\ii\mu_1k\nu_{\beta1} & -\mu_2\Gamma_2 & -2\ii\mu_2k\nu_{\beta2}\\[2pt]
 -2\ii\mu_1k\nu_{\alpha1} & -\mu_1\Gamma_1 & -2\ii\mu_2k\nu_{\alpha2} & \mu_2\Gamma_2
\end{bmatrix}}_{\mathbf M_1}
\begin{bmatrix}A_1\\ B_1\\ A_2\\ B_2\end{bmatrix}
=
-\begin{bmatrix}
 \ii k\,P^0-\nu_{\beta1}S^0\\
 \nu_{\alpha1}P^0+\ii k\,S^0\\
 \mu_1\Gamma_1P^0+2\ii\mu_1k\nu_{\beta1}S^0\\
 2\ii\mu_1k\nu_{\alpha1}P^0-\mu_1\Gamma_1S^0
\end{bmatrix}.
\label{eq:case1}
\end{equation}
Rows: $u_x$, $u_z$, $\sigma_{zz}$, $\sigma_{xz}$. The dispersion relation
$\Delta_1(k)\equiv\det\mathbf M_1=0$ with $k>\max(k_{\beta1},k_{\beta2})$ real is the
\emph{Stoneley} equation. A true (non-leaky) Stoneley root exists only for a limited range of
contrasts (roughly $\beta_1\approx\beta_2$); otherwise the root is leaky (see
Section~\ref{sec:leaky}).

\subsection{Case 2: solid--fluid, source in the solid}

Conditions at $z=0$: (i) $u_z^{s}=u_z^{f}$, (ii) $\sigma_{zz}^{s}=-p$, (iii) $\sigma_{xz}^{s}=0$.
Unknowns $\mathbf x=[A_1,\,B_1,\,A_f]^{\mathsf T}$:
\begin{equation}
\underbrace{\begin{bmatrix}
 -\nu_{\alpha} & \ii k & -\nu_f\\[2pt]
 \mu\Gamma & -2\ii\mu k\nu_{\beta} & \rho_f\omega^2\\[2pt]
 -2\ii\mu k\nu_{\alpha} & -\mu\Gamma & 0
\end{bmatrix}}_{\mathbf M_2}
\begin{bmatrix}A_1\\ B_1\\ A_f\end{bmatrix}
=
-\begin{bmatrix}
 \nu_{\alpha}P^0+\ii k\,S^0\\
 \mu\Gamma P^0+2\ii\mu k\nu_{\beta}S^0\\
 2\ii\mu k\nu_{\alpha}P^0-\mu\Gamma S^0
\end{bmatrix}
\label{eq:case2}
\end{equation}
(solid subscript $1$ dropped). Expanding along the third column,
\begin{equation}
  \Delta_2(k)\equiv\det\mathbf M_2
  =\mu^2\Bigl[\nu_f\,R(k)+\frac{\rho_f}{\rho}\,k_\beta^4\,\nu_\alpha\Bigr],
  \label{eq:scholte}
\end{equation}
and $\Delta_2=0$ with $k>\max(k_\beta,k_f)$ real is the \emph{Scholte} equation. A true Scholte root
always exists.

\subsection{Case 3: solid--fluid, source in the fluid}

Same interface conditions. Unknowns $\mathbf x=[A_1,\,B_1,\,A_f]^{\mathsf T}$ (transmitted P, SV into the
solid; reflected P in the fluid). The incident fluid field $P_f^0e^{-\nu_fz}$ has
$u_z=-\nu_fP_f^0$ and $\sigma_{zz}=-\rho_f\omega^2P_f^0$, which gives
\begin{equation}
\mathbf M_2
\begin{bmatrix}A_1\\ B_1\\ A_f\end{bmatrix}
=
\begin{bmatrix}
 -\nu_f\,P_f^0\\
 -\rho_f\omega^2\,P_f^0\\
 0
\end{bmatrix}.
\label{eq:case3}
\end{equation}
The matrix is the same as in Case~2, so the poles (Scholte, leaky Rayleigh) are the same.
Only their excitation differs. This is a statement of source--receiver reciprocity.

%======================================================================
\section{Degeneration to the free surface (Lamb/Rayleigh)}
%======================================================================

\paragraph{Matrix level.} With $\rho_f=0$, $\mathbf M_2$ becomes block lower-triangular:
\begin{equation}
\mathbf M_2\big|_{\rho_f=0}=
\begin{bmatrix}
 -\nu_\alpha & \ii k & -\nu_f\\
 \mu\Gamma & -2\ii\mu k\nu_\beta & 0\\
 -2\ii\mu k\nu_\alpha & -\mu\Gamma & 0
\end{bmatrix}
=\begin{bmatrix}
 \ast & -\nu_f\\
 \mathbf M_{\mathrm{FS}} & \mathbf 0
\end{bmatrix},
\qquad
\mathbf M_{\mathrm{FS}}=
\begin{bmatrix}
 \mu\Gamma & -2\ii\mu k\nu_\beta\\
 -2\ii\mu k\nu_\alpha & -\mu\Gamma
\end{bmatrix}.
\end{equation}
Rows 2--3 decouple into the traction-free system
$\mathbf M_{\mathrm{FS}}[A_1,B_1]^{\mathsf T}=\mathbf S_{2:3}$
($\sigma_{zz}=\sigma_{xz}=0$ at $z=0$). Row~1 then only gives the passive fluid motion $A_f$,
which carries no energy because its flux is proportional to $\rho_f$
(Section~\ref{sec:energy}). Also,
\begin{equation}
  \det\mathbf M_{\mathrm{FS}}=-\mu^2\bigl[\Gamma^2-4k^2\nu_\alpha\nu_\beta\bigr]=-\mu^2R(k),
\end{equation}
which is consistent with \eqref{eq:scholte}: $\Delta_2\to\mu^2\nu_fR(k)$ as $\rho_f\to0$. The zeros of
$\Delta_2$ in the solid sheet therefore tend to the Rayleigh root $R(k_R)=0$, $k_R=\omega/c_R$.

\paragraph{Pole level (weak loading).} Let $\varepsilon=\rho_f/\rho$. Expanding
\eqref{eq:scholte} about $k_R$:
\begin{equation}
  k_p=k_R+\delta k,\qquad
  \delta k=-\,\varepsilon\,\frac{k_\beta^4\,\nu_\alpha(k_R)}{\nu_f(k_R)\,R'(k_R)}+O(\varepsilon^2).
  \label{eq:dk}
\end{equation}
\begin{itemize}
  \item \emph{Fast fluid}, $\alpha_f>c_R$ ($k_f<k_R$): $\nu_f(k_R)$ is real, so $\delta k$ is real.
  The pole is a true (real) Scholte pole near $k_R$. As $\varepsilon\to0$ it tends continuously to
  the Rayleigh pole, and so does its share of the energy.
  \item \emph{Slow fluid}, $\alpha_f<c_R$ ($k_f>k_R$, e.g.\ water on rock, or $\alpha_f\to0$):
  $\nu_f(k_R)$ is imaginary, $\delta k$ is complex, and the pole is a \emph{leaky Rayleigh} pole,
  $k_L=a+\ii b$ with $b>0$, on the sheet $\RE\nu_f<0$. It radiates into the fluid over the leakage
  length $1/(2b)$. For $\alpha_f\to0$, $\nu_f(k_R)\simeq-\ii\omega/\alpha_f$ and
  $b\propto\rho_f\alpha_f$. The controlling small parameter is the fluid impedance
  $Z_f=\rho_f\alpha_f$.
  \item The real Scholte pole $k_{\mathrm{Sch}}>k_f$ goes to $\infty$ as $\alpha_f\to0$, and its
  energy vanishes.
  \item \emph{Non-uniform limit.} In the slow-fluid route, the energy partition at infinite range
  sends the leaky-Rayleigh energy to the fluid P wave for every $Z_f>0$. The free-surface partition
  is recovered only for ranges $r\ll1/(2b)$. On the real axis inside the fluid window, the leaky pole
  appears as a Lorentzian peak of half-width $b$ in $|D_{\mathrm P_f}|^2$ and in
  $\IM\mathcal S$. The quadrature must resolve that peak with nodes graded at $a\pm b\,10^{j}$;
  otherwise it is missed \emph{on both sides} of the energy balance, and the balance still closes.
  The pole is located with real-axis Newton steps $k\leftarrow a-\Delta(a)/\Delta'(a)$,
  $a\leftarrow\RE k$, which are exact to $O(b^2)$ because $\Delta$ is analytic inside the window.
  We report the peak energy as \emph{quasi-trapped}. For $\rho_f=\alpha_f=10^{-6}$ this gives
  $b\approx1.8\times10^{-13}$ ($\omega=\beta=1$), and the Miller--Pursey partition is recovered
  to $10^{-2}\,\%$.
  \item Case~3 in this limit: $\mathbf S\propto(\nu_f,\rho_f\omega^2,0)^{\mathsf T}P_f^0$, and
  $A_1,B_1\to0$. The fluid sees a rigid wall and no energy enters the solid.
\end{itemize}

%======================================================================
\section{Energy partition}
\label{sec:energy}
%======================================================================

\subsection{Flux density}

For $e^{-\ii\omega t}$ fields, the time-averaged energy flux is
$\langle\mathcal P_j\rangle=-\tfrac12\RE(\sigma_{ij}\dot u_i^{*})$, i.e.
\begin{equation}
  \boxed{\ \langle\mathcal P_j\rangle=\frac{\omega}{2}\,\IM\bigl(\sigma_{ij}\,u_i^{*}\bigr)\ }
\end{equation}
(power per unit length of the line source). Energy is quadratic, so the partition must be computed
at infinity, where the far fields of the different wave types separate and their cross terms vanish.
Squaring the individual terms of a contour decomposition at finite range does not give a valid
partition.

\subsection{Contour decomposition of the wavefield}

For $x>0$ the inverse transform is closed in the upper half $k$-plane. Using Cramer's rule,
$\mathbf x(k)=\operatorname{adj}\mathbf M\,\mathbf S/\Delta(k)$, and
\begin{equation}
  \mathbf u(x,z)=
  \underbrace{\ii\sum_{p}\Res_{k=k_p}\bigl[\mathbf u(k,z)\bigr]e^{\ii k_px}}_{\text{interface waves, }O(1)\text{ in }x}
  \;+\;
  \underbrace{\frac{1}{2\pi}\sum_{m}\int_{\mathcal B_m}\mathbf u(k,z)\,e^{\ii kx}\,\dd k}_{\text{body/head waves, }O(r^{-1/2})},
  \label{eq:contour}
\end{equation}
where $\mathcal B_m$ are the branch cuts from $k_m$ ($m=\alpha_1,\beta_1,\alpha_2,\beta_2$ or $f$).
The sum includes only poles on the physical sheet \eqref{eq:branch}. For a lossless medium these are
real, with $k_p>\max_m k_m$. As $\omega\to\omega+\ii0$ the pole $+k_p$ lies at $k_p+\ii0$
(outgoing to $x>0$) and $-k_p$ lies at $-k_p-\ii0$.

\paragraph{Residue (numerically stable form).} Let $\mathbf r,\bm\ell$ be the right and left null
vectors of $\mathbf M(k_p)$. Then
\begin{equation}
  \Res_{k=k_p}\mathbf M^{-1}=\frac{\mathbf r\,\bm\ell^{\mathsf T}}{\bm\ell^{\mathsf T}\mathbf M'(k_p)\,\mathbf r},
  \qquad
  \Res_{k=k_p}\mathbf x=\frac{\mathbf r\,\bigl(\bm\ell^{\mathsf T}\mathbf S(k_p)\bigr)}{\bm\ell^{\mathsf T}\mathbf M'(k_p)\,\mathbf r},
  \label{eq:res}
\end{equation}
which is equivalent to $\operatorname{adj}\mathbf M(k_p)\mathbf S(k_p)/\Delta'(k_p)$. $\mathbf M'$ can be
evaluated analytically or by a centred difference. The incident field has no pole, so only the
scattered amplitudes contribute.

\subsection{Interface-wave power (pole)}

For $x>0$ the interface-wave field has amplitudes
$\mathbf x^{\mathrm{IW}}=\ii\,\Res_{k_p}\mathbf x$ and phase $e^{\ii k_px}$. At the pole all $\nu_m$
are real and positive. The horizontal power through a vertical plane is
\begin{equation}
  W_{\mathrm{IW}}^{+}=\frac{\omega}{2}\int_{-\infty}^{\infty}
  \IM\bigl(\sigma_{xx}u_x^{*}+\sigma_{xz}u_z^{*}\bigr)\,\dd z ,
  \label{eq:WIW}
\end{equation}
with fields from the table in Section~1, $k\to k_p$, and $a,b\to$ the entries of
$\mathbf x^{\mathrm{IW}}$. In the fluid, use $\sigma_{xx}=-\rho_f\omega^2\phi_f$ and
$\sigma_{xz}=0$. Every product is of the form $e^{-(\nu_m+\nu_n)|z|}$, so the integral can be done
analytically: $\int_0^\infty e^{-(\nu_m+\nu_n)z}\dd z=(\nu_m+\nu_n)^{-1}$. For $x<0$,
$W_{\mathrm{IW}}^{-}$ follows in the same way from $-k_p$ with amplitude $-\ii\Res_{-k_p}\mathbf x$. For a
symmetric source (explosion, vertical force), $W_{\mathrm{IW}}^{-}=W_{\mathrm{IW}}^{+}$, and
\begin{equation}
  W_{\mathrm{IW}}=W_{\mathrm{IW}}^{+}+W_{\mathrm{IW}}^{-}.
\end{equation}

\subsection{Body-wave power (real axis / branch cuts)}

The body waves are the branch-cut integrals in \eqref{eq:contour}. Their power is obtained without
evaluating those integrals. Take the vertical flux through the horizontal planes $z=\pm Z$, which lie
beyond the source and outside the evanescent interface-wave layer, and apply Parseval's theorem
$\int|f|^2\dd x=\tfrac1{2\pi}\int|\tilde f|^2\dd k$:
\begin{equation}
  \int_{-\infty}^{\infty}\langle\mathcal P_z\rangle\,\dd x
  =\frac{\omega}{2}\,\frac{1}{2\pi}\int_{-\infty}^{\infty}
  \IM\bigl(\sigma_{xz}u_x^{*}+\sigma_{zz}u_z^{*}\bigr)\,\dd k .
\end{equation}
For a single plane wave of amplitude $D$ (P, SV, or fluid P, travelling away from the interface),
substituting the table of Section~1 gives
\begin{equation}
  \boxed{\ \mathcal E_m(k)=\tfrac12\,\rho\,\omega^3\,\RE\bigl(\ii\nu_m\bigr)
  =\begin{cases}\tfrac12\rho\omega^3\sqrt{k_m^2-k^2}, & |k|<k_m\\ 0,& |k|>k_m\end{cases}\ }
  \label{eq:Em}
\end{equation}
per unit $|D|^2$. For propagating waves the up/down and P/S cross terms vanish, so the fluxes add.
Hence
\begin{equation}
  W_{\mathrm{body}}=\frac{1}{2\pi}\sum_{m}\int_{-k_m}^{k_m}\mathcal E_m(k)\,\bigl|D_m(k)\bigr|^2\,\dd k ,
  \label{eq:Wbody}
\end{equation}
with the outgoing amplitudes $D_m$ listed below. Direct and reflected waves are added coherently
because they cannot be separated in the far field.
\begin{center}
\renewcommand{\arraystretch}{1.35}
\begin{tabular}{@{}llll@{}}
\toprule
Case & Region & Wave $m$ & $D_m$ (with $\rho$ of that medium in $\mathcal E_m$)\\
\midrule
1 & $z>h$ (solid 1) & P$_1$, S$_1$ & $P_d^0+A_1$,\quad $S_d^0+B_1$\\
1 & $z<0$ (solid 2) & P$_2$, S$_2$ & $A_2$,\quad $B_2$\\
2 & $z>h$ (solid)   & P, S       & $P_d^0+A_1$,\quad $S_d^0+B_1$\\
2 & $z<0$ (fluid)   & P$_f$      & $A_f$\\
3 & $z<-h$ (fluid)  & P$_f$      & $U_f^0+A_f$\\
3 & $z>0$ (solid)   & P, S       & $A_1$,\quad $B_1$\\
\bottomrule
\end{tabular}
\end{center}
Poles lie outside every window $|k|<k_m$, so \eqref{eq:Wbody} has no pole singularities. It does
have integrable inverse-square-root singularities at \emph{every} branch point $k_n\le k_m$ inside the
window. The window endpoint contributes the $1/\nu_m$ Weyl factor. An interior point $k_n<k_m$ gives,
for example, the grazing-P $\to$ SV conversion at $k=k_{\alpha1}$ inside the S window, which is the
head-wave contribution. Split each window at all branch points $0<k_{n_1}<\dots<k_m$, and on each
segment $[k_a,k_b]$ substitute
\begin{equation}
  k^2=k_a^2+(k_b^2-k_a^2)\sin^2 t,\qquad t\in[0,\tfrac{\pi}{2}]
  \qquad(k=k_b\sin t\ \text{if}\ k_a=0),
  \label{eq:subst}
\end{equation}
which regularises both endpoints; then apply Gauss--Legendre quadrature in $t$. Plain midpoint
or trapezoid rules in $k$ or $\theta$ converge erratically, with errors of order $10^{-3}$. With
$k=k_m\sin\theta$, \eqref{eq:Wbody} also gives the angular radiation pattern $\dd W/\dd\theta$.

\subsection{Partition percentages and energy-balance check}

\begin{equation}
  W_{\mathrm{tot}}=W_{\mathrm{body}}+W_{\mathrm{IW}},\qquad
  \eta_m=100\,\frac{W_m}{W_{\mathrm{tot}}}\ [\%],\qquad
  \eta_{\mathrm{IW}}=100\,\frac{W_{\mathrm{IW}}}{W_{\mathrm{tot}}}\ [\%].
\end{equation}
\paragraph{Independent check (vertical force).} The power delivered by the source is
$W_{\mathrm{src}}=\tfrac12\RE(F^{*}\dot u_z)=\tfrac{\omega}{2}\IM\bigl(F^{*}u_z(0,h)\bigr)$.
By the Sokhotski--Plemelj formula, $(k-k_p\mp\ii0)^{-1}=\mathrm{PV}\,(k-k_p)^{-1}\pm\ii\pi\delta(k-k_p)$,
which gives
\begin{equation}
  W_{\mathrm{src}}=\frac{\omega}{2}\left[\frac{1}{2\pi}\,\mathrm{PV}\!\!\int_{-\infty}^{\infty}
  \IM\bigl(F^{*}u_z(k,h)\bigr)\dd k
  \;+\;\RE\Bigl(F^{*}\Res_{k=k_p}u_z(k,h)\Bigr)\right]
  \label{eq:Wsrc}
\end{equation}
(both poles $\pm k_p$ included). The second term equals $W_{\mathrm{IW}}$ and the first equals
$W_{\mathrm{body}}$. The code must satisfy $W_{\mathrm{src}}=W_{\mathrm{tot}}$ to quadrature accuracy.
For the explosive source the analogous check uses the injected volume and the pressure at the source.

\paragraph{Verified reference values.} These values come from a numerical evaluation of this
formulation in the free-surface limit, with a vertical force and a Poisson solid
($\alpha=\sqrt3\,\beta$, $c_R=0.9194\,\beta$). The energy balance \eqref{eq:Wsrc} closes to a relative
error of about $10^{-8}$:
\begin{center}
\begin{tabular}{@{}ccccc@{}}
\toprule
$k_\beta h$ & P [\%] & S [\%] & Rayleigh [\%] & comment\\
\midrule
$0$   & 21.94 & 29.13 & 48.93 & same at every $\omega$ (no other length scale)\\
$0.4$ & 20.34 & 26.89 & 52.77 & \\
$1.2$ & 17.48 & 20.68 & 61.83 & \\
$3$   &  3.18 & 25.91 & 70.91 & \\
$9$   & 12.52 & 87.19 &  0.29 & Rayleigh excitation $\sim e^{-2\nu_\beta(k_R)h}$\\
\bottomrule
\end{tabular}
\end{center}
The dependence on depth is not monotonic because $W_{\mathrm{src}}$ changes with $h$ through
interference with the surface reflection.

%======================================================================
\section{Leaky poles}
\label{sec:leaky}
%======================================================================

Complex poles on non-physical sheets (the leaky Rayleigh wave in Case~2/3 when $\alpha_f>c_R$, and
leaky Stoneley waves in Case~1) radiate continuously into body waves. They are \emph{not} trapped
energy. Their energy is already counted in \eqref{eq:Wbody}, as a peak of $|D_m(k)|^2$ near
$\RE k_p$. Report them, if needed, as a sub-partition of $W_{\mathrm{body}}$. They must not be added
to $W_{\mathrm{IW}}$.

%======================================================================
\section{Three-dimensional formulation: point force and general moment tensor}
\label{sec:3d}
%======================================================================

\subsection{Transform and local frame}

The source is at $(0,0,z_s)$. We use the double Fourier transform
\begin{equation}
  \tilde f(k_x,k_y,z)=\iint f\,e^{-\ii(k_xx+k_yy)}\,\dd x\,\dd y,\qquad
  f(x,y,z)=\frac{1}{4\pi^2}\iint \tilde f\,e^{\ii(k_xx+k_yy)}\,\dd k_x\,\dd k_y ,
\end{equation}
with polar wavenumbers $k_x=k\cos\varphi$, $k_y=k\sin\varphi$ ($k\ge0$) and the local frame
\begin{equation}
  \hat{\mathbf l}=(\cos\varphi,\sin\varphi,0)^{\mathsf T},\qquad
  \hat{\mathbf t}=(-\sin\varphi,\cos\varphi,0)^{\mathsf T},\qquad \hat{\mathbf z}.
\end{equation}
For fixed $(k,\varphi)$ every field varies horizontally only as $e^{\ii k\,\hat{\mathbf l}\cdot\mathbf x}$.
The problem therefore splits exactly into:
\begin{itemize}
  \item a \textbf{P--SV problem} in the vertical plane $(\hat{\mathbf l},\hat{\mathbf z})$. It is the 2D
  problem of Sections~1--5 with $u_x\to u_L=\mathbf u\cdot\hat{\mathbf l}$ and
  $\sigma_{xz}\to\sigma_{Lz}$, and \emph{the matrices $\mathbf M_1,\mathbf M_2$ are unchanged};
  \item an \textbf{SH problem} for $u_T=\mathbf u\cdot\hat{\mathbf t}$. With
  $u_T=C\,e^{s\nu_\beta z}$: $\sigma_{zT}=s\,\mu\,\nu_\beta\,C$ and $\sigma_{LT}=\ii k\mu\,C$.
\end{itemize}
The 3D Weyl identity has the same kernel as in 2D:
\begin{equation}
  g_m=\frac{e^{\ii k_mR}}{4\pi R}
  \quad\Longleftrightarrow\quad
  \hat g_m(k,z)=\frac{e^{-\nu_m|z-z_s|}}{2\nu_m}.
\end{equation}

\subsection{Source representation}

The whole-space Green's tensor ($u_i=G_{ij}F_j$) is
\begin{equation}
  G_{ij}=\frac{1}{\rho\omega^2}\Bigl[k_\beta^2\,\delta_{ij}\,g_\beta+\partial_i\partial_j\bigl(g_\beta-g_\alpha\bigr)\Bigr].
\end{equation}
A point force $\mathbf F$ (any direction) radiates $u_i=G_{ij}F_j$. A moment tensor $\mathbf M$ (any
$3\times3$ matrix; symmetric in the seismic case) radiates
$u_i=M_{pq}\,\partial G_{ip}/\partial\xi_q=-M_{pq}\,\partial_qG_{ip}$. For each wave type define the
(bilinear) wave vector
\begin{equation}
  \bm\kappa_m^{\epsilon}=\bigl(\ii k\cos\varphi,\ \ii k\sin\varphi,\ -\epsilon\,\nu_m\bigr)^{\mathsf T},
  \qquad \epsilon=\sgn(z-z_s),\qquad
  \bm\kappa_m\!\cdot\bm\kappa_m=-k_m^2 .
\end{equation}
For $z\neq z_s$ the direct field is $\tilde{\mathbf u}^0=c_P\,\bm\kappa_\alpha\,\hat g_\alpha+\mathbf w\,\hat g_\beta$, where
\begin{align}
  \text{force:}\quad
  c_P&=-\frac{\bm\kappa_\alpha\!\cdot\mathbf F}{\rho\omega^2}, &
  \mathbf w&=\frac{1}{\rho\omega^2}\Bigl[k_\beta^2\,\mathbf F+\bm\kappa_\beta\,(\bm\kappa_\beta\!\cdot\mathbf F)\Bigr],
  \label{eq:srcF}\\
  \text{moment:}\quad
  c_P&=\frac{\bm\kappa_\alpha^{\mathsf T}\mathbf M\,\bm\kappa_\alpha}{\rho\omega^2}, &
  \mathbf w&=-\frac{1}{\rho\omega^2}\Bigl[k_\beta^2\,\mathbf M\bm\kappa_\beta+\bm\kappa_\beta\,(\bm\kappa_\beta^{\mathsf T}\mathbf M\,\bm\kappa_\beta)\Bigr],
  \label{eq:srcM}
\end{align}
with every $\bm\kappa$ evaluated at the same $\epsilon$. Here $\mathbf w$ is transverse
($\bm\kappa_\beta\!\cdot\mathbf w=0$). Projecting onto the potentials of Section~1 and onto the SH
displacement gives
\begin{equation}
  \boxed{\ \Pi^{\epsilon}=\frac{c_P}{2\nu_\alpha},\qquad
  \Sigma^{\epsilon}=\frac{w_z}{2\ii k\,\nu_\beta},\qquad
  H^{\epsilon}=\frac{\mathbf w\cdot\hat{\mathbf t}}{2\nu_\beta}\ }
  \label{eq:PSH}
\end{equation}
so that $\phi^0=\Pi^\epsilon e^{-\nu_\alpha|z-z_s|}$, $\psi^0=\Sigma^\epsilon e^{-\nu_\beta|z-z_s|}$ and
$u_T^0=H^\epsilon e^{-\nu_\beta|z-z_s|}$. For a source in the solid ($z_s=h$),
\begin{equation}
  (P^0,S^0,C^0)=(\Pi^{-}e^{-\nu_\alpha h},\ \Sigma^{-}e^{-\nu_\beta h},\ H^{-}e^{-\nu_\beta h}),\qquad
  (P^0_d,S^0_d,C^0_d)=(\Pi^{+}e^{\nu_\alpha h},\ \Sigma^{+}e^{\nu_\beta h},\ H^{+}e^{\nu_\beta h}).
\end{equation}
$P^0,S^0$ enter $\mathbf S$ in \eqref{eq:case1}--\eqref{eq:case2} unchanged, and $C^0$ enters the SH
system below.

\paragraph{Consistency.} (i) With $\mathbf F=F\hat{\mathbf z}$ these formulas reproduce the 2D
vertical-force table. (ii) For $\mathbf M=M_0\mathbf I$ they give $\mathbf w=\mathbf 0$ and
$c_P=-M_0/(\lambda+2\mu)$, i.e.\ the explosion with $\Phi_0=-M_0/(\lambda+2\mu)$.
(iii) \emph{Source in the fluid} (Case~3, $z_s=-h$): only the P terms survive, with
$\rho\to\rho_f$ and $\alpha\to\alpha_f$:
$c_P=-\bm\kappa_f\!\cdot\mathbf F/(\rho_f\omega^2)$ or
$\bm\kappa_f^{\mathsf T}\mathbf M\bm\kappa_f/(\rho_f\omega^2)$, $P_f^0=\Pi_f^{+}e^{-\nu_fh}$, and
$U_f^0=\Pi_f^{-}e^{\nu_fh}$.

\paragraph{Azimuthal structure.} $\Pi,\Sigma,H$ are trigonometric polynomials in $\varphi$ of degree
$\le1$ for a force and $\le2$ for a moment tensor. These are the azimuthal orders $m=0,\pm1,\pm2$ of
the cylindrical theory. Hence $|D|^2$ has degree $\le4$, and the uniform trapezoidal rule in
$\varphi$ with $N_\varphi\ge5$ nodes is \emph{exact}. Cartesian displacement components have degree
$\le3$.

\subsection{SH boundary conditions}

The unknowns are $C_1$ (reflected SH, $e^{-\nu_{\beta1}z}$) and $C_2$ (transmitted SH, $e^{+\nu_{\beta2}z}$).
\begin{align}
  \text{Case 1:}&\quad
  \begin{bmatrix}1 & -1\\ -\mu_1\nu_{\beta1} & -\mu_2\nu_{\beta2}\end{bmatrix}
  \begin{bmatrix}C_1\\ C_2\end{bmatrix}=-\begin{bmatrix}C^0\\ \mu_1\nu_{\beta1}C^0\end{bmatrix},
  \qquad \Delta_{\mathrm{SH}}=-(\mu_1\nu_{\beta1}+\mu_2\nu_{\beta2}),\\
  \text{Case 2 and free surface:}&\quad \sigma_{zT}=0\ \Rightarrow\ C_1=C^0\quad\text{(total reflection)},\\
  \text{Case 3:}&\quad C_1\equiv0\quad\text{(a fluid source radiates no SH)}.
\end{align}
$\Delta_{\mathrm{SH}}\neq0$ on the physical sheet, so two half-spaces carry \emph{no} SH interface wave.
Love waves need a layer. SH energy is therefore always body-wave energy.

\subsection{Energy partition in 3D}

By Parseval, $\iint|f|^2\dd x\,\dd y=\frac{1}{4\pi^2}\iint|\tilde f|^2\dd^2k$. The flux densities
\eqref{eq:Em} still apply per unit $|{\rm potential}|^2$ for P and SV. For SH, per unit
$|C|^2$, the flux density is
\begin{equation}
  \mathcal E_{\mathrm{SH}}(k)=\tfrac12\,\omega\,\mu\,\RE(\ii\nu_\beta).
\end{equation}

\paragraph{Body waves.}
\begin{equation}
  \boxed{\ W_m=\frac{1}{4\pi^2}\int_0^{2\pi}\!\!\dd\varphi\int_0^{k_m}\mathcal E_m(k)\,
  \bigl|D_m(k,\varphi)\bigr|^2\,k\,\dd k\ },
  \qquad m\in\{\mathrm P_1,\mathrm{SV}_1,\mathrm{SH}_1,\mathrm P_2,\mathrm{SV}_2,\mathrm{SH}_2,\mathrm P_f\},
  \label{eq:Wbody3}
\end{equation}
with $D_m$ from the table of Section~\ref{sec:energy} plus
$D_{\mathrm{SH}_1}=C_d^0+C_1$ (Cases~1,~2) and $D_{\mathrm{SH}_2}=C_2$ (Case~1).
Integrate in $k$ with \eqref{eq:subst} and in $\varphi$ with the exact trapezoidal rule.

\paragraph{Interface wave.} The pole $k_p$ does not depend on $\varphi$ because the medium is isotropic.
A residue in $k$ followed by stationary phase in $\varphi$ gives the far field
\begin{equation}
  \mathbf u^{\mathrm{IW}}(r,\theta,z)\simeq
  \ii\,\sqrt{\frac{k_p}{2\pi r}}\;e^{\ii(k_pr-\pi/4)}\;
  \Res_{k=k_p}\tilde{\mathbf u}(k,\varphi=\theta,z),
\end{equation}
and the power through a cylinder $r\to\infty$ is
\begin{equation}
  \boxed{\ W_{\mathrm{IW}}=\frac{k_p}{2\pi}\int_0^{2\pi}\mathcal W^{+}_{\mathrm{2D}}(\varphi)\,\dd\varphi\ },
  \label{eq:WIW3}
\end{equation}
where $\mathcal W^{+}_{\mathrm{2D}}(\varphi)$ is \eqref{eq:WIW} evaluated in the $(\hat{\mathbf l},\hat{\mathbf z})$
plane with amplitudes $\Res_{k_p}\mathbf x(k,\varphi)$. The $u_\theta$ terms decay as $r^{-3/2}$ and
drop out.

\paragraph{Source work (independent check, force or moment).}
\begin{equation}
  W_{\mathrm{src}}=\frac{\omega}{2}\,\IM\,\mathcal S[\mathbf u],\qquad
  \mathcal S[\mathbf u]=
  \begin{cases}\mathbf F^{*}\!\cdot\mathbf u(\mathbf x_s) & \text{force},\\
  M^{*}_{pq}\,\partial_qu_p(\mathbf x_s) & \text{moment}.\end{cases}
\end{equation}
Split $\mathbf u$ into the direct field and the scattered field. The direct field contributes the
whole-space power $W_0$, which can be computed with \eqref{eq:Wbody3} applied to the direct up- and
down-going amplitudes. For a force it is also available analytically:
$W_0=\omega^2|\mathbf F|^2(\alpha^{-3}+2\beta^{-3})/(24\pi\rho)$. For the scattered field (regular
at $\mathbf x_s$), $\partial_q\to\kappa_q$ of each wave. For the lower solid the polarisation and wave
vectors are
\begin{align*}
  &\text{P: } A\,e^{-\nu_\alpha z}(\ii k\hat{\mathbf l}-\nu_\alpha\hat{\mathbf z}),\ \
  \bm\kappa=\ii k\hat{\mathbf l}-\nu_\alpha\hat{\mathbf z};\qquad
  \text{SV: } B\,e^{-\nu_\beta z}(\nu_\beta\hat{\mathbf l}+\ii k\hat{\mathbf z}),\ \
  \bm\kappa=\ii k\hat{\mathbf l}-\nu_\beta\hat{\mathbf z};\\
  &\text{SH: } C\,e^{-\nu_\beta z}\,\hat{\mathbf t},\ \ \bm\kappa=\ii k\hat{\mathbf l}-\nu_\beta\hat{\mathbf z};
  \qquad \mathcal S_{\mathbf M}=\textstyle\sum_{\text{waves}}\mathbf v^{\mathsf T}\mathbf M^{*}\bm\kappa .
\end{align*}
Sokhotski--Plemelj on the radial integral (a single pole at $k_p$, since $\varphi$ covers all
directions) gives
\begin{equation}
  W_{\mathrm{src}}=W_0
  +\frac{\omega}{8\pi^2}\,\mathrm{PV}\!\!\int_0^{2\pi}\!\!\!\int_0^{\infty}\IM\,\mathcal S[\tilde{\mathbf u}^{\mathrm{sc}}]\,k\,\dd k\,\dd\varphi
  +\underbrace{\frac{\omega k_p}{8\pi}\int_0^{2\pi}\RE\,\mathcal S\bigl[\Res_{k_p}\tilde{\mathbf u}^{\mathrm{sc}}\bigr]\dd\varphi}_{=\,W_{\mathrm{IW}}}.
  \label{eq:Wsrc3}
\end{equation}
The last term is a second, independent evaluation of \eqref{eq:WIW3}.

\subsection{Verification (3D, numerical evaluation of this section)}

Poisson solid, $\rho=\beta=\omega=1$, free-surface limit. The general moment tensor used is
$\mathbf M=\left[\begin{smallmatrix}1&0.4&-0.7\\0.4&-0.3&0.5\\-0.7&0.5&0.6\end{smallmatrix}\right]$ and the
oblique force is $\mathbf F=(0.3,-0.5,0.8)$.
\begin{center}
\small
\renewcommand{\arraystretch}{1.2}
\begin{tabular}{@{}lll@{}}
\toprule
Test & Result & Reference\\
\midrule
Whole space, oblique $\mathbf F$: $W_0$ & $0.02849672$ & analytic $0.02849671$\\
Whole space, force: $W_S/W_P$ & $10.3923$ & $2(\alpha/\beta)^3=10.3923$\\
Whole space, double couple: $W_S/W_P$ & $23.3827$ & $\tfrac32(\alpha/\beta)^5=23.3827$\\
Surface vertical force: P / S / R [\%] & $6.88$ / $25.75$ / $67.36$ & Miller--Pursey: $6.9$ / $25.8$ / $67.4$\\
$\mathbf M$, $k_\beta h=0.5$: P/SV/SH/R [\%] & $2.69$ / $8.48$ / $62.02$ / $26.81$ & balance err.\ $1.8\times10^{-8}$\\
$\mathbf F$, $k_\beta h=0.5$: P/SV/SH/R [\%] & $6.27$ / $21.63$ / $18.64$ / $53.46$ & balance err.\ $3.6\times10^{-8}$\\
\bottomrule
\end{tabular}
\end{center}
In both buried cases, \eqref{eq:WIW3} (flux) and the last term of \eqref{eq:Wsrc3} (source side)
agree to all printed digits.

\subsection{Wavefield synthesis with the discrete wavenumber method}

Sample $\varphi$ uniformly ($N_\varphi\ge8$). Take the FFT in $\varphi$ of each Cartesian component,
$\tilde{\mathbf u}(k,\varphi,z)=\sum_n\tilde{\mathbf u}_n(k,z)e^{\ii n\varphi}$ with $|n|\le3$. The
Jacobi--Anger identity $\int_0^{2\pi}e^{\ii n\varphi}e^{\ii kr\cos(\varphi-\theta)}\dd\varphi=2\pi\ii^nJ_n(kr)e^{\ii n\theta}$
then gives
\begin{equation}
  \mathbf u(r,\theta,z)=\frac{1}{2\pi}\sum_{n=-3}^{3}\ii^n e^{\ii n\theta}
  \int_0^\infty\tilde{\mathbf u}_n(k,z)\,J_n(kr)\,k\,\dd k
  \;\approx\;\frac{\Delta k}{2\pi}\sum_{n}\ii^n e^{\ii n\theta}\sum_{j\ge1}\tilde{\mathbf u}_n(k_j,z)J_n(k_jr)\,k_j,
\end{equation}
with $k_j=j\,\Delta k$, $\Delta k=2\pi/L$ (Bouchon's DWN). $L$ must exceed $r+c_{\max}T$ so that the
periodic images arrive after the time window $T$.

\emph{Complex frequency.} With the $e^{-\ii\omega t}$ convention the contour is shifted \emph{upward},
$\omega_c=\omega+\ii\epsilon$, and
$u(t)=e^{\epsilon t}\,\frac{1}{2\pi}\int U(\omega_c)e^{-\ii\omega t}\dd\omega$. In MATLAB this synthesis is
\texttt{fft} (kernel $e^{-\ii\omega t}$), not \texttt{ifft}. A typical choice is $\epsilon\approx\pi/T$ to
$2\pi/T$.

\emph{Energy partition is evaluated at real $\omega$.} The damping $\epsilon$ dissipates energy and
moves $k_p$ off the real axis, which blurs the separation between body and interface waves. The
energy module therefore removes the pole analytically through \eqref{eq:res}, and the body windows
contain no poles. Optional cross-check: $W_{\mathrm{src}}(\omega+\ii\epsilon)\to W_{\mathrm{tot}}$ as
$\epsilon\to0$, because the Lorentzian peak at $k_p$ integrates to the pole term.

%======================================================================
\section{Remarks for the MATLAB implementation}
%======================================================================

\begin{enumerate}
  \item Scale the stress rows by $1/\mu_1$ (and $\mathbf S$ accordingly) and use
  $\bar k=k/k_{\beta1}$. This keeps $\mathbf M$ well conditioned.
  \item Find the root $k_p$ by scanning $k>\max_mk_m$ for sign changes of $\Delta(k)$, which is
  real there up to a constant phase because all $\nu$ are real, and then refine with \texttt{fzero}.
  The Scholte root lies beyond $k_f$ and can be far out when $\alpha_f$ is small.
  \item Plane-wave check: for $|k|<k_{\alpha1}$ with incident P only,
  $\sum_m\mathcal E_m|x_m|^2/(\mathcal E_{\alpha1}|P^0|^2)=1$ (energy reflection and transmission
  coefficients).
  \item Broadband partition: weight $W_m(\omega)$ by $|\hat s(\omega)|^2$ of the source wavelet and
  integrate over $\omega$.
  \item 3D: see Section~\ref{sec:3d}. The same $\mathbf M_1,\mathbf M_2$ are evaluated for each
  $(k,\varphi)$, together with a scalar or $2\times2$ SH system.
\end{enumerate}

\end{document}
